\part{Unidad 2}

\chapter{Cómputo en la nube}

\section{Aprovisionamiento de recursos informáticos On-premise}

En el pasado, el aprovisionamiento de recursos informáticos On-premise (en las instalaciones de la empresa) representaba el concurso de diferentes proveedores de bienes y servicios, como eran los representantes de ventas, ingenieros de preventa, fabricantes de los equipos, fabricante del sistema operativo, fabricante de la base de datos, agentes aduanales, transportistas, instaladores del site, proveedor de energía, proveedor de comunicaciones, instaladores del hardware, instaladores del software, entre otros.

Entonces, el aprovisionamiento de recursos informáticos era un proceso complejo y tardado, el cual culminaba con el sistema en producción.

Después había que reaprovisionar cuando crecían las necesidades de la empresa.

\section{Cómputo como servicio}

El concepto de cómputo en la nube (cloud computing) está basado en las ideas del profesor John McCarthy de la Universidad de Stanford.

En un discurso en 1961, predijo que la tecnología de tiempo compartido (la posibilidad de que múltiples usuarios y procesos utilicen una misma computadora) podría llevar en el futuro, a ofrecer las capacidades de cómputo y las aplicaciones como servicios públicos.

Esta idea se concretaría en lo que actualmente conocemos como cómputo en la nube, al ofrecer servicios virtualizados a nivel masivo en Internet.

Por otra parte, en 1984 John Burdette Gage acuñó el slogan The network is the computer para Sun Microsystems, el cual significa que la computadora de escritorio es una ``ventana`` a la red, y mediante el software apropiado, es posible transferir tareas de cómputo a otras computadoras.

A la posibilidad de trasferir tareas de cómputo a otras computadoras se le conoce como computation offloading.

\section{Fábricas de información}

En 2006 aparece el artículo ``The Information Factories`` en la prestigiosa revista WIRED [1].

En este artículo se presenta la idea de ``fábricas de información`` (Information Factories).

El autor, George Gilder, argumenta que la tecnología de la información evoluciona hacia un modelo de procesamiento y almacenamiento de datos centralizados en grandes centros de datos.

Estos centros de datos estarían equipados con grandes servidores y redes de alta velocidad.

Entonces, en lugar de depender de computadoras personales, los usuarios accederán a servicios informáticos a través de redes de comunicación utilizando dispositivos más simples, como serían las computadoras portátiles y los dispositivos móviles como tabletas y teléfonos celulares.

> ``The desktop is dead. Welcome to the Internet cloud, where massive facilities across the globe will store all the data you'll ever use``, The Information Factories, George Gilder.

Este modelo de informática centralizada permitiría contar con capacidades ``ilimitadas`` de procesamiento y almacenamiento, lo cual proporcionará niveles de escalabilidad sin precedentes.

Si bien George Gilder no introduce el concepto de cómputo en la nube como tal, su artículo explora las implicaciones sociales y económicas de este modelo de computación basado en la idea de ``fábrica de información``.

\section{Proveedores de servicios en la nube}

\subsection{Amazon}

En 2006 Amazon lanzó la plataforma Amazon Web Services (AWS) incluyendo los servicios Elastic Computer (EC2) y Storage Cloud (S3).

AWS popularizó el modelo de Infraestructura como Servicio (IaaS, Infrastructure as a Service), el cual permite a las empresas alquilar recursos de cómputo y almacenamiento en lugar de invertir en hardware propio.

\subsection{Google}

En 2008 Google lanzó Google App Engine, inicialmente como una plataforma para desarrollar y alojar aplicaciones web.

Posteriormente, Google lanza el servicio Google Cloud Storage en 2010 y en 2013 el servicio Google Compute Engine (GCE). Con estos servicios Google entra en el mercado de infraestructura como Servicio (IaaS).

La marca Google Cloud Platform (GCP) se lanzó oficialmente en 2011.

\subsection{Microsoft}

En 2010 Microsoft lanzó la plataforma Windows Azure (ahora conocida como Microsoft Azure), con servicios para desarrolladores al nivel de plataforma (PaaS), entre otros:

\begin{itemize}

\item Windows Azure Compute para la ejecución de aplicaciones en la nube utilizando máquinas virtuales administradas por Microsoft.

\item Windows Azure Storage para el almacenamiento de datos estructurados y no estructurados.

\item Windows Azure SQL Database el cual es SQL Server en la nube.

\end{itemize}

En 2012 Microsoft lanzó los servicios:

\begin{enumerate}

\item Azure Virtual Machines, con el cual Microsoft entra en el mercado de IaaS.

\item Azure Virtual Network el cual permite crear redes virtuales en la nube.

\item Azure Storage IaaS el cual ofrece almacenamiento para las máquinas virtuales.

\end{enumerate}

En 2015 Microsoft lanza el servicio Azure Container Service el cual permite ejecutar contenedores Docker en Azure.

\section{Definición de cómputo en la nube}

El Instituto Nacional de Estándares y Tecnología (NIST) de Estados Unidos, define el cómputo en la nube de la siguiente manera [2]:

> ``Cloud computing is a model for enabling ubiquitous, convenient, on-demand network access to a shared pool of configurable computing resources (e.g., networks, servers, storage, applications, and services) that can be rapidly provisioned and released with minimal management effort or service provider interaction.``

\section{Las cinco características esenciales del cómputo en la nube}

\subsection{1. Autoservicio bajo demanda}

Los usuarios pueden unilateralmente aprovisionar automáticamente capacidades en la infraestructura del proveedor del servicio de nube, tales como tiempo de servidor, almacenamiento y red, de acuerdo a sus necesidades sin intervención humana.

\subsection{2. Acceso amplio por medio de la red}

Las capacidades que ofrece el proveedor de nube están disponibles a través de la red, y son accedidas mediante mecanismos estándares utilizando clientes heterogéneos como teléfonos celulares, tabletas, computadoras de escritorio, etc.

\subsection{3. Recursos compartidos}

Los recursos de cómputo del proveedor de nube son compartidos por múltiples usuarios usando un modelo “multiinquilino” (multitenant). Los recursos físicos y virtuales (procesamiento, almacenamiento, memoria, ancho de banda, etc.) son asignados y reasignados dinámicamente de acuerdo a las necesidades de los usuarios.

Los recursos asignados son transparentes en términos de la ubicación, se manera que el usuario no tiene que saber qué recurso físico está utilizando. No obstante, el usuario puede determinar en qué región o datacenter creará los recursos virtuales.

\subsection{4. Elasticidad}

Las capacidades pueden ser aprovisionadas y des-aprovisionadas elásticamente (manualmente o automáticamente), permitiendo al usuario escalar hacia arriba o hacia abajo, horizontalmente o verticalmente, bajo demanda. La escalabilidad hacia arriba con frecuencia parece ser ilimitada.

\subsection{5. Medición del servicio}

Los recursos de nube pueden ser monitoreados con la finalidad de controlar y optimizar automáticamente su utilización. Este monitoreo se realiza a diferentes niveles de abstracción de acuerdo al tipo de servicio (almacenamiento, procesamiento, ancho de banda, cuentas de usuario, etc.).

\section{Los tres modelos de servicio del cómputo en la nube}

\subsection{1. Infraestructura como servicio}

En el nivel de Infraestructura como servicio o Infrastructure as a Service (IaaS), el usuario puede aprovisionar procesamiento (máquinas virtuales), almacenamiento (discos virtuales), redes y otros servicios que le permiten ejecutar cualquier software.

En este nivel el usuario tiene el control de sistema operativo, el almacenamiento, algunos componentes de la red, y las aplicaciones que ejecutan sobre el sistema operativo.

\subsection{2. Plataforma como servicio}

En el nivel de plataforma como servicio o Platform as a Service (PaaS), el usuario puede aprovisionar servicios de instancias de bases de datos, respaldos, ejecución de servicio web, plataformas de desarrollo y pruebas, y otros servicios sin la necesidad de aprovisionar o controlar la infraestructura subyacente (servicios de IaaS como serían la máquina virtual, la red o el almacenamiento).

\subsection{3. Software como servicio}

En el nivel de software como servicio o Software as a Service (SaaS), el usuario puede utilizar las aplicaciones publicadas por el proveedor de nube, las cuales ejecutan sobre la infraestructura de nube. El usuario no tiene que aprovisionar servicios de PaaS o de IaaS.

Los servicios al nivel de SaaS se acceden mediante una aplicación cliente, p.e. un navegador (web browser). Los usuarios acceden a las aplicaciones de SaaS mediante un usuario y contraseña o mediante un par de claves.

\section{Los cuatro modelos de despliegue del cómputo en la nube}

\subsection{1. Nube privada}

Los servicios de nube son aprovisionados para el uso exclusivo de una sola organización.

La organización puede ser propietaria de la infraestructura, de un tercero o de una combinación de ambas.

De igual forma, la administración y operación puede ser realizada por la misma organización, por un tercero o una combinación de ambas.

En cuanto a la ubicación, la infraestructura puede estar on-premise (en las instalaciones de la organización), o bien, off-premise (fuera de las instalaciones de la organización).

\subsection{2. Nube comunitaria}

Los servicios de nube son aprovisionados para el uso de una comunidad de usuarios u organizaciones con objetivos compartidos (p.e. empresas, instituciones de gobierno, etc.).

Al igual que la nube privada, la propiedad, administración y operación de la infraestructura puede ser de la comunidad, un tercero o una combinación de ambos.

La infraestructura puede estar ubicada en instalaciones de la comunidad, un tercero o una combinación de ambas.

\subsection{3. Nube pública}

Los servicios de nube son aprovisionados para el público en general.

La infraestructura está ubicada en las instalaciones del proveedor de los servicios de nube (Amazon, Microsoft, Google, etc.).

\subsection{4. Nube híbrida}

La infraestructura de nube es una composición de dos o más modelos de despliegue (nube privada, nube comunitaria, nube pública).

La comunicación entre las redes de las diferentes nubes, se realiza mediante tecnologías propietarias o estándares que permitan implementar canales de comunicación seguros (como son las VPN).

\chapter{Principios del cómputo en la nube}

\section{Servicios virtuales a nivel masivo}

El cómputo en la nube es un modelo de arquitectura basado en una infraestructura de cómputo ofrecida como servicios virtuales a nivel masivo.

El concepto clave en el cómputo en la nube es el ``servicio``, así, se ofrece:

\begin{itemize}

\item IaaS (Infrastructure as a Service) Servicio de creación de infraestructura virtual y física.

\item PaaS (Platform as a Service) Servicios de DBMS, plataformas de desarrollo y pruebas.

\item SaaS (Software as a Service) Servicio de aplicaciones de software.

\item Otros servicios con la terminación ``as a Service``:

\item DaaS (Data as a Service) Servicio de acceso a datos.

\item DRaaS (Disaster Recovery as a Service) Servicio de recuperación de máquinas virtuales ante desastres.

\end{itemize}

\section{Escalabilidad en la nube}

La escalabilidad es la capacidad de un sistema en la nube de aumentar o disminuir los recursos aprovisionados (CPU virtuales y memoria de una máquina virtual, disco de sistema operativo, disco de datos.

La escalabilidad vertical se refiere al crecimiento o descrecimiento de las capacidades de una computadora.

La escalabilidad horizontal se refiere a la creación o eliminación de máquinas virtuales.

\section{Disponibilidad en la nube}

La disponibilidad es la capacidad de un sistema en la nube para estar accesible (disponible) en todo momento, de manera que los usuarios puedan acceder el sistema.

La disponibilidad se mide en porcentaje que representa el tiempo durante el cual el sistema está disponible (tiempo de actividad o uptime).

La disponibilidad es: \textbf{100 x el tiempo de actividad / tiempo total}.

Por ejemplo, si el sistema está disponible 364 días de 365 días, entonces la disponibilidad es del 99.7260\%

Los proveedores de nube establecen la disponibilidad de los servicios en un SLA (Service Level Agreement) o Acuerdo de nivel de servicio.

\section{Confiabilidad en la nube}

La confiabilidad es la capacidad de un sistema en la nube para funcionar durante largos períodos de tiempo sin interrupciones.

Para garantizar la confiabilidad, los proveedores de nube implementan diversas estrategias, por ejemplo:

\begin{itemize}

\item \textbf{Geo-Redundant Storage.} La replicación de los respaldos en regiones a cientos de kilómetros de la región actual.

\item \textbf{Zona de disponibilidad.} Uno o más datacenters equipados con su propia alimentación, refrigeración y red.

\item \textbf{Dominios de error.} Las máquinas virtuales que pertenecen a un mismo dominio de error comparten la misma fuente de alimentación y red física.

\item \textbf{Dominios de actualización.} Las máquinas virtuales dentro de un mismo dominio de actualización se reiniciarán simultáneamente durante el mantenimiento planeado. Solo se reiniciará un dominio de actualización a la vez.

\item \textbf{Balanceo de carga.} Un balanceador de carga permite escalar un sistema, y al mismo tiempo evita puntos únicos de falla aumentando la disponibilidad.

\item \textbf{Site Recovery.} Permite replicar y recuperar máquinas virtuales en caso de falla o desastre.

\item \textbf{Respaldos.} Copia de seguridad de máquinas virtuales y archivos, permitiendo la recuperación de los datos en caso de ransomware, pérdida o corrupción de los datos

\item \textbf{Alta disponibilidad en DBMS.} Crea una réplica de la base de datos en una zona de disponibilidad diferente.

\end{itemize}

\section{Transparencia}

La transparencia es la capacidad de los usuarios para acceder los recursos de nube sin la necesidad de conocer los detalles de la infraestructura de nube.

Podemos dividir la transparencia en siete categorías:

\begin{itemize}

\item \textbf{Transparencia en el acceso a los datos.} Los usuarios y aplicaciones acceden a los datos en la nube de manera estandarizada.

\item \textbf{Transparencia de ubicación.} Los usuarios de nube acceden a los recursos independientemente de su localización física.

\item \textbf{Transparencia de migración.} Los usuarios de nube no son afectados por la migración de los datos y de los procesos que pueda realizar la plataforma.

\item \textbf{Transparencia de reubicación.} La transparencia de reubicación se refiere a la capacidad de la nube de cambiar la ubicación de un recurso mientras está en uso, sin que el usuario que accede el recurso se vea afectado.

\item \textbf{Transparencia de replicación.} La transparencia de replicación es la capacidad de la nube de ocultar la existencia de recursos replicados.

\item \textbf{Transparencia de concurrencia.} La transparencia de concurrencia se refiere a la capacidad de la nube de ocultar el hecho de que varios usuarios y procesos comparten los diferentes recursos de manera concurrente.

\item \textbf{Transparencia ante fallas.} La transparencia ante fallas es la capacidad de la nube de ocultar una falla.

\end{itemize}

Por ejemplo, si un sistema que se encuentra totalmente replicado, cuando el sistema principal falla entonces el usuario accederá de manera transparente a la réplica del sistema.

Más adelante en el curso veremos cómo replicar un sistema completo en la nube utilizando un administrador de tráfico de red y cómo implementar sistemas tolerantes a fallas en la nube mediante balance de carga.

\section{Rendimiento}

El rendimiento es la eficiencia con la que un sistema en la nube utiliza los recursos de nube, incluyendo la capacidad de procesamiento, almacenamiento, ancho de banda, entre otros recursos.

\section{Desempeño}

El desempeño se refiere a la velocidad con la ejecuta un sistema en la nube.

El desempeño depende de las técnicas de programación y del rendimiento del sistema. El objetivo es mantener el sistema funcionando de manera consistente incluso bajo cargas de trabajo pesadas.

\section{La elasticidad en la nube}

Debido a que el cómputo en la nube está basado fundamentalmente en la virtualización de los recursos informáticos, este modelo de arquitectura ofrece una ventaja única, la posibilidad de hacer crecer y decrecer los recursos aprovisionados.

Supongamos un servicio de streaming bajo demanda, como es el caso de Netflix. En este tipo de servicio la demanda crece los fines de semana y decrece los días entre semana.

Si el proveedor del servicio no aprovisiona los recursos suficientes para atender la demanda del fin de semana, entonces muchos usuarios se quedarán sin servicio.

Por otra parte, si el proveedor del servicio aprovisiona los recursos necesarios para atender a sus usuarios el fin de semana, estos recursos estarán subutilizados los días entre semana, lo cual resulta en pérdidas económicas.

Sin lugar a dudas, el éxito que han alcanzado las empresas proveedoras de streaming bajo demanda, se debe a que su modelo de negocio está basado en la posibilidad que les ofrece la nube para crecer y decrecer los recursos aprovisionados, a esta característica de la nube se le llama elasticidad.

El cómputo elástico es la habilidad de hacer crecer y decrecer rápidamente la capacidad de cómputo (CPUs), la memoria y el almacenamiento para adaptarse a la demanda.

Para implementar el cómputo elástico se utilizan herramientas de monitoreo, las cuales aprovisionan y des-aprovisionan recursos conforme son necesarios, sin detener la operación.

Sin embargo, hay proveedores de nube que afirman que no existe tal cosa como ``nube privada``, ya que el tema de elasticidad se ve acotado en un equipo ``privado`` debido a la limitada escalabilidad.

En cambio, en la nube pública la escalabilidad es casi ilimitada, ya que si se agotan los recursos de un data center, sin ningún problema se puede escalar a otro data center, ya sea del mismo proveedor o de otros proveedores.

Una característica importante de la nube (pública) es ``pagar por lo que se usa``, esto significa que solo se paga por los recursos aprovisionados.

En el caso de un centro de cómputo tradicional (también es el caso de la ``nube privada``) la empresa paga por todo el equipo, lo use a toda su capacidad o no.

Consideremos la siguiente gráfica:

\subsection{Escenario On-premise}

Al inicio los requerimientos informáticos de la empresa son pocos, sin embargo ésta debe hacer una gran inversión inicial ya que debe adquirir el equipo que le permita operar durante un periodo de tiempo determinado.

Para saber cuál será la inversión inicial  se deberá hacer una planeación de capacidad.

Suponiendo que la empresa crece, entonces en determinado momento deberá realizar la actualización de su equipo informático.

Entonces se deberá hacer una planeación de capacidad para establecer el tamaño de la inversión, estimando un tiempo determinado de operación.

No obstante las previsiones, la empresa podría decrecer, por tanto los requerimientos informáticos también decrecen. El costo de la actualización no se recupera, incluso, en poco tiempo el equipo se hará obsoleto.

\subsection{Escenario nube}

Ahora supongamos que la empresa utiliza servicios en la nube. Entonces la inversión inicial es mínima, ya que la empresa contratará los servicios indispensables para operar. Conforme la empresa crece, la elasticidad de la nube permite adecuar el tamaño de la infraestructura informática de acuerdo a las necesidades.

Si la empresa decrece, también decrecen sus necesidades y el costo de la infraestructura informática. En la gráfica, las barras corresponden al costo de los servicios de nube. Podemos ver que este costo se ajusta a las necesidades de la empresa.

Podemos concluir que al utilizar los servicios de nube se optimiza la relación costo/beneficio, ya que solo pagamos lo que realmente usamos.
